% !TEX root = ../../../main.tex
% Appendix entry: preprocessing.bandpass -- "Bandpass filter" (Appendix A).
\entryfield{Algorithm}
\code{_run} in \code{Adapters/preprocessing_bandpass.py} first refuses
$\mathit{low\_hz} \geq \mathit{high\_hz}$ with a \code{ValueError}, then calls the factory
\code{make_bandpass_filter} in \code{Working/Detection/analysis/freq_analysis.py} and applies
the function it returns to the span. With sampling rate $f_s$, order $N$ and edges
$f_\mathrm{lo} < f_\mathrm{hi}$:
\begin{enumerate}
\item Normalise the edges by the Nyquist frequency $f_s/2$:
  $W = \bigl(2f_\mathrm{lo}/f_s,\ 2f_\mathrm{hi}/f_s\bigr)$.
\item Design a digital Butterworth band-pass filter \citep{butterworth1930theory} with
  \code{scipy.signal.butter(N, W, btype="band")} \citep{virtanen2020scipy}, returned as
  transfer-function coefficients $(b, a)$, each of length $2N+1$ (a band-pass of order $N$ is a
  filter of order $2N$).
\item Run it forwards and then backwards over the span with
  \code{scipy.signal.filtfilt(b, a, x)}, after casting $x$ to floating point.
\end{enumerate}
Because the filter runs both ways, the result has no phase shift (an event stays at the sample
where it happened), and its magnitude response is the square of the single-pass Butterworth
response:
\[ |H_\mathrm{ff}(f)| = |H_\mathrm{B}(f)|^2 , \]
so each band edge is attenuated to one half in amplitude ($-6$\,dB), not to $1/\sqrt2$
($-3$\,dB) as in the single-pass design, and the roll-off outside the band is twice as steep.

\code{filtfilt} is called with SciPy's defaults: before filtering it extends the span at each
end by an odd reflection about the end sample ($2x_0 - x_{k}$ at the start and the mirror at
the end) of length $3\max(\mathrm{len}(a), \mathrm{len}(b)) = 3(2N+1)$ samples (27 at the
default $N = 4$), starts the filter from the steady state for that end value
(\code{lfilter_zi}), and trims the extension off afterwards. The output has the span's length.

\entryfield{Assumptions}
\begin{itemize}
\item \emph{Sampling rate.} The edges are in hertz and converted with the span's own $f_s$, so
  the same settings select the same physical band on a 1\,Hz and a 10\,Hz recording. Both edges
  must lie strictly between 0 and $f_s/2$: the registry enforces only a minimum of
  $10^{-6}$\,Hz, and SciPy's \code{butter} rejects a normalised edge at or above 1 with a
  \code{ValueError}. The default band list in \code{Working/run_groups.py}
  (\code{DEFAULT_BANDS}) stops its third band at 0.45\,Hz for this reason.
\item \emph{Units.} The block is linear and not unit-aware (its \code{_run} takes no
  \code{units} argument); the output is in the recording's stored unit, and scaling the input
  scales the output by the same factor.
\item \emph{Missing values.} Nothing removes or bridges NaNs. In an infinite-impulse-response
  filter a NaN enters the filter's state and every later output becomes NaN; the backward pass
  then carries it to every earlier output, so one NaN anywhere in the span makes the whole
  output NaN. \unverified{No test checks this; it follows from how \code{lfilter} propagates
  state.}
\item \emph{Span length.} \code{filtfilt} refuses a span no longer than its padding,
  $3(2N+1)$ samples (27 at $N = 4$, 63 at $N = 10$), with a \code{ValueError}.
\item \emph{Mean.} A band-pass removes the zero-frequency part, so the output sits near zero
  whatever the recording's baseline.
\end{itemize}

\entryfield{Outputs}
A \texttt{Signal} with the filtered samples and the input's $f_s$. Nothing is put in
\code{meta}. Like every block's output it is held by the chain's step cache; the block writes
nothing of its own.

\entryfield{Cost}
Two passes of a recursive filter of order $2N$: $O(nN)$ time and $O(n)$ memory, vectorised in
SciPy. Fast on a whole channel; no estimate is registered and it is not meant for the cluster.

\entryfield{Failure modes}
\begin{itemize}
\item \emph{An edge at or above Nyquist} (for example 0.5\,Hz on a 1\,Hz recording) stops the
  step with SciPy's error about critical frequencies; an edge \emph{close} to Nyquist is
  accepted.
\item \emph{One NaN blanks the span} (above).
\item \emph{Ringing and pre-ringing.} A step or a sharp drop passed through a band-pass gives
  an oscillating response; because the filter runs both ways the oscillation appears
  symmetrically, \emph{before} the event as well as after it. Detectors downstream can mark
  these lobes as events of the opposite sign, and depths measured on the filtered trace are not
  the depths on the raw trace.
\item \emph{Edge transients.} The odd reflection pads only $3(2N+1)$ samples, far fewer than
  the filter's impulse response is long when the lower edge is a small fraction of $f_s$ (its
  response lasts on the order of $1/f_\mathrm{lo}$ seconds: 100\,s at the 0.01\,Hz default).
  Expect distortion over roughly that length at each end of the span.
  \unverified{The length of the distorted region has not been measured.}
\item \emph{Numerical accuracy at very low edges.} The filter is designed and applied in
  transfer-function form $(b, a)$, not as second-order sections (\code{sosfiltfilt}). For a
  band whose edges are tiny fractions of $f_s$ (the 0.001--0.01\,Hz default band on a 1\,Hz
  recording normalises to $W = (0.002, 0.02)$, at order 8 overall) this form is known to lose
  precision. \unverified{Whether it materially distorts the output at the shipped defaults has
  not been checked.}
\end{itemize}

\entryfield{Validation}
No test in the repository checks the filter's frequency response or its output values.
SciPy's \code{butter} and \code{filtfilt} are standard and widely tested
\citep{virtanen2020scipy}; what is untested here is the wrapping (edge normalisation, the
choice of transfer-function form, the defaults).

\entryfield{References}
\citet{butterworth1930theory} for the filter family; \citet{virtanen2020scipy} for the
implementation.

\entryfield{Tests}
\code{tests/test_t06_adapter_remap.py} checks the declared types and that the output is a
\texttt{Signal} of the input's length and rate. \code{tests/test_chain_validation.py} uses it
to check chain type-validation. \code{tests/test_band_scope.py} checks that a band run inserts
exactly this block with the band's edges (and that the surrogate run is band-passed the same
way), and that an upside-down band in Settings is refused. \code{tests/test_artifacts.py}
checks figure file names. None of them checks which frequencies pass.
